% !PRINTED 136
\weberpage{136}

Aus (3) und (4) folgt für die Doppelverhältnisse:
\weq{5}{\frac{(\xi_2\eta_3)(\xi_1\eta_4)}{(\xi_1\eta_3)(\xi_2\eta_4)} = \frac{\vartheta^4_{10}}{\vartheta^4_{00}},}
\weq{6}{\frac{(\xi_1\eta_2)(\xi_3\eta_4)}{(\xi_1\eta_3)(\xi_2\eta_4)} = \frac{\vartheta^4_{01}}{\vartheta^4_{00}}.}
Wenn wir nun zwei der Gleichungen (1), etwa die beiden ersten,
differentiieren, so folgt:
\[
\begin{aligned}
2\vartheta_{01}(u)\vartheta'_{01}(u)\,du &= \eta_1\,d\xi - \xi_1\,d\eta = (\eta_1\,d\xi),\\
2\vartheta_{11}(u)\vartheta'_{11}(u)\,du &= \eta_2\,d\xi - \xi_2\,d\eta = (\eta_2\,d\xi),
\end{aligned}
\]
und daraus mit Benutzung von (1)
\[
\begin{aligned}
2\vartheta_{01}(u)\vartheta_{11}(u)\bigl[\vartheta'_{01}(u)\vartheta_{11}(u) - \vartheta'_{11}(u)\vartheta_{01}(u)\bigr]\,du \\
= (\eta_1\,d\xi)\vartheta^2_{11}(u) - (\eta_2\,d\xi)\vartheta^2_{01}(u) \\
= (\eta_1\,d\xi)(\xi\eta_2) - (\eta_2\,d\xi)(\xi\eta_1) = (\xi_2\eta_1)(\xi\,d\eta).
\end{aligned}
\]
Den letzten Ausdruck kann man ohne Rechnung dadurch ableiten, daß man
den vorletzten als lineare Funktion von \(d\xi\), \(d\eta\) betrachtet, die für
\(d\xi : d\eta = \xi : \eta\) verschwindet und daher durch \((\xi\,d\eta)\) teilbar ist. Den
Quotienten \((\xi_2\eta_1)\) erhält man, wenn man \(d\xi : d\eta = \xi_2 : \eta_2\) setzt. Mit
Benutzung von \S\,23, (6) erhält man dann
\weq{7}{2\pi\vartheta^2_{01}\vartheta_{00}(u)\vartheta_{11}(u)\vartheta_{10}(u)\vartheta_{01}(u)\,du = (\xi_1\eta_2)(\xi\,d\eta).}
Führt man hierin nach (3) (erste und letzte Formel)
\weq{8}{(\xi_1\eta_2)\vartheta^2_{00} = \sqrt{(\xi_1\eta_3)(\xi_2\eta_4)}\ \vartheta^2_{01}}
ein, und setzt für die \(\vartheta\)-Funktionen die Ausdrücke (1), so folgt
schließlich
\weq{9}{2\pi\vartheta^2_{00}\,du = \frac{\sqrt{(\xi_1\eta_3)(\xi_2\eta_4)}\,(\xi\,d\eta)}{\sqrt{(\xi\eta_1)(\xi\eta_2)(\xi\eta_3)(\xi\eta_4)}},}
wodurch \(du\) als \gesperrt{elliptisches Differential erster Gattung} in homogenen
Variablen [\S\,1, (5)] dargestellt ist.

Es ist durch (1) das Verhältnis \(\xi : \eta\) als \gesperrt{doppeltperiodische Funktion}
von \(u\) bestimmt. Desgleichen sind aber auch die Verhältnisse der
Quadratwurzeln
\weq{10}{\sqrt{(\xi\eta_1)},\qquad \sqrt{(\xi\eta_2)},\qquad \sqrt{(\xi\eta_3)},\qquad \sqrt{(\xi\eta_4)}}
als eindeutige doppeltperiodische Funktionen erklärt, und das Vorzeichen
der Quadratwurzeln in (9) ist hierdurch und durch (8) ebenfalls eindeutig
bestimmt.
